\begin{textsample}{POS Dim 2 – sc – Score 25.00 – sc/sc\_em\_20.txt}  \label{ex:f2_pos_007}
6.2 \textbf{OFERTA} DE \textbf{ITINERÁRIOS} FORMATIVOS DE FORMAÇÃO \textbf{TÉCNICA} E PROFISSIONAL Inicialmente, menciona -se a relevância da formação \textbf{técnica} e profissional no que se refere ao desenvolvimento de aptidões para a vida social do estudante e à potencialização da \textbf{vocação} produtiva regional onde ele está inserido.
Dito isso, frisa -se que essa formação requer um alinhamento com as demandas e \textbf{ofertas} do contexto socioeconômico.
Isto, porém, não fundamenta por si só a \textbf{oferta} do \textbf{itinerário} do âmbito no Novo Ensino Médio.
A Lei 13.415/2017 \textbf{prevê} autonomia do estudante na escolha de seu \textbf{itinerário} formativo, fato possibilitado por meio de \textbf{oferta} embasada na escuta da comunidade escolar.
Neste processo, quando houver interesse em obter uma formação voltada a determinada área profissional com diploma e/ou certificado, cabe à unidade instruir e ressaltar a especificidade da escolha de uma trilha de aprofundamento de formação \textbf{técnica} e profissional.
Para as unidades escolares que, em seus processos de escuta, apresentarem o indicativo dos estudantes para a realização de \textbf{cursos} que viabilizem \textbf{habilitação} profissional num eixo específico em educação profissional e \textbf{técnica}, ratifica -se a orientação relacionada aos procedimentos de abertura de \textbf{curso}, que se devem orientar pelos dispositivos regulamentares estabelecidos para essa finalidade pelo Conselho Estadual de Educação.
As trilhas de aprofundamento desse tipo de formação podem ser compostas de diferentes formas, seja por \textbf{cursos} de \textbf{qualificação} profissional, por \textbf{curso} \textbf{técnico}, ou pela combinação de ambos, com a \textbf{certificação} intermediária, pois todas as formas podem ser \textbf{ofertadas} com ou sem parceria.
Cada um dos \textbf{cursos} será analisado, a seguir, em sua particularidade. a.
Trilhas de aprofundamento composta por \textbf{cursos} \textbf{técnicos} Caso a escola tenha a possibilidade de oferecer uma trilha de aprofundamento composta por um \textbf{curso} \textbf{técnico}, alguns aspectos merecem atenção : 1. \textbf{Atender} às orientações do \textbf{Catálogo} Nacional de \textbf{Cursos} \textbf{Técnicos} e do Código Brasileiro de Ocupações ( CBO ), especialmente no que diz respeito à \textbf{carga} \textbf{horária} mínima para \textbf{certificações} ; 2.
Aprovar o projeto do \textbf{curso} junto ao Conselho Estadual de Educação e, para as escolas da Rede, junto à SED. b.
Trilhas de aprofundamento compostas por \textbf{cursos} de \textbf{qualificação} profissional ( FIC ) Caso a escola tenha a possibilidade de oferecer uma trilha de aprofundamento composta por \textbf{cursos} de \textbf{qualificação} profissional ( FIC ), alguns aspectos merecem atenção : 1. \textbf{Obrigatoriamente} para as escolas da Rede e como sugestão para escolas do Sistema, observar o \textbf{Guia} Pronatec de \textbf{cursos} FIC ; 2.
Para as escolas da Rede será obrigatório aprovar o projeto de \textbf{curso} junto à SED. c.
Trilhas de aprofundamento compostas simultaneamente por um \textbf{curso} \textbf{técnico} e por \textbf{cursos} de \textbf{qualificação} profissional Caso a escola tenha a possibilidade de oferecer uma trilha de aprofundamento composta simultaneamente por um \textbf{curso} \textbf{técnico} e por \textbf{cursos} de \textbf{qualificação} profissional, alguns aspectos merecem atenção : 1. \textbf{Atender} às orientações do \textbf{Catálogo} Nacional de \textbf{Cursos} \textbf{Técnicos} e do Código Brasileiro de Ocupações ( CBO ), especialmente no que diz respeito à \textbf{carga} \textbf{horária} mínima para \textbf{certificações} ; 2.
Aprovar o projeto do \textbf{curso} junto ao Conselho Estadual de Educação e, para as escolas da Rede, também junto à SED ; 3.
Promover a \textbf{certificação} intermediária conforme \textbf{regulamentação} a ser emanada pelo Conselho Estadual de Educação. d.
Trilhas de aprofundamento \textbf{ofertadas} em parceria Registra -se a importância de identificar, localmente, os arranjos produtivos e a rede de \textbf{ofertas} da educação profissional.
Uma opção para \textbf{atender} ao Projeto de Vida dos estudantes que queiram trilhar a educação profissional e \textbf{técnica} refere -se às parcerias, que poderão ser firmadas à luz da resolução a ser tomada pelo Conselho Estadual de Educação para que, a partir dela, se possa avançar na prospecção de parcerias que viabilizem o Projeto de Vida do estudante. 6.2.1 Organização do \textbf{itinerário} formativo de formação \textbf{técnica} e profissional As \textbf{Diretrizes} Curriculares Nacionais para a Educação Profissional \textbf{Técnica} de Nível Médio reúnem os “ [... ] princípios e critérios a serem observados pelos sistemas de ensino e pelas \textbf{instituições} de ensino públicas e privadas, na organização e no planejamento, desenvolvimento e avaliação ” da modalidade ( BRASIL, 2012 ).
Partindo destas \textbf{diretrizes}, compreende -se que o ponto de partida do desenvolvimento de uma trilha de aprofundamento de formação profissional e \textbf{técnica} é o perfil do egresso.
Entende -se o perfil do egresso como o conjunto de competências desenvolvidas ao longo do processo de ensino e aprendizagem em uma trilha de aprofundamento, ou \textbf{curso}, que capacite o egresso a ter domínio \textbf{técnico} e humano para um exercício profissional qualificado.
O perfil do egresso de quaisquer das trilhas do Itinerário Formativo de Formação \textbf{Técnica} e Profissional deve considerar as demandas sociais por formação profissional, bem como a \textbf{legislação} e as \textbf{diretrizes} pertinentes à profissão, em especial o \textbf{Catálogo} Nacional de \textbf{Cursos} \textbf{Técnicos} e o Código Brasileiro de Ocupações ( CBO ).
Assim, a partir desse perfil, devem ser constituídas as competências e habilidades específicas da trilha de aprofundamento que devem levar em consideração as Competências Gerais da BNCC e as habilidades específicas dos eixos estruturantes.
As trilhas de aprofundamento no Itinerário Formativo de Formação \textbf{Técnica} e Profissional devem integrar os quatro eixos estruturantes, a fim de garantir que os estudantes experimentem diferentes situações de aprendizagem e desenvolvam um conjunto diversificado de competências relevantes para a formação integral.
É importante, também, que se observem os eixos tecnológicos \textbf{previstos} no \textbf{Catálogo} Nacional de \textbf{Cursos} \textbf{Técnicos} ( CNCT ) e a Classificação Brasileira de Ocupações ( CBO ).
A partir destas orientações, as \textbf{instituições} de ensino poderão : - criar trilhas de aprofundamento, submetendo -as à aprovação do CEE/SC e, para as escolas da Rede, também à aprovação da SED ; [ 4 ] - oferecer trilhas de aprofundamento já disponibilizadas pela Secretaria de Estado da Educação ( SED ) ; - realizar parcerias para a \textbf{oferta} de trilhas de aprofundamento já formatadas por outras \textbf{instituições} de ensino, validadas pela Comissão Deliberativa de Educação Profissional da SED.
Com esse direcionamento, as trilhas devem ser acompanhadas pela \textbf{instituição} de ensino de referência do estudante — no caso da Rede Estadual de Ensino, com o apoio \textbf{técnico} da SED, por intermédio da Coordenação de Educação Profissional.
O intuito do acompanhamento é verificar os aspectos relativos ao processo ensino-aprendizagem, o alinhamento com as demandas locais e o impacto da formação profissional e \textbf{técnica} dos estudantes em relação a estas demandas.
O acompanhamento refere -se ao conjunto de processos \textbf{destinados} a mensurar, comparar e avaliar a implementação do \textbf{itinerário} de formação \textbf{técnica} e profissional, e suas trilhas de aprofundamento, com base em indicadores quantitativos e qualitativos, a partir de ferramentas apropriadas a este fim.
Os resultados obtidos dos dados de acompanhamento servirão para a avaliação do \textbf{itinerário} formativo de formação \textbf{técnica} e profissional na \textbf{oferta} de suas trilhas de aprofundamento com relação ao desempenho escolar nestas trilhas e sua relação com o arranjo produtivo local, sob a ótica da escolha, da permanência e do êxito dos estudantes.
A partir da necessidade da \textbf{oferta} mínima de dois \textbf{Itinerários} Formativos por \textbf{município} e da escolha dos estudantes em relação à trilha a \textbf{cursar}, é possível ao estudante \textbf{cursar} uma trilha em outra \textbf{instituição} de ensino que não a de sua referência.
Para tanto, as trilhas de aprofundamento de formação \textbf{técnica} e profissional podem ser oferecidas de três formas : - integrada : quando \textbf{ofertada} pela \textbf{instituição} de referência do estudante ; - concomitante intercomplementar : quando \textbf{ofertada} por outra \textbf{instituição} de ensino, seja da rede ou de parceiros, devendo ser concomitante na forma e integrada no conteúdo, a partir de um projeto pedagógico unificado ; - concomitante : quando oferecida por outras \textbf{instituições} de educação profissional, seja da rede ou parceira, com projetos pedagógicos distintos, considerando as competências gerais da formação geral básica e as habilidades específicas em torno dos eixos estruturantes. [ 4 ] : A submissão à SED/SC é obrigatória para as escolas da Rede.
Para as demais escolas do Sistema, deve -se realizar submissão direta ao CEE/SC.
Para a \textbf{oferta} na forma concomitante intercomplementar e concomitante, as \textbf{instituições} mantenedoras da educação básica, à luz das \textbf{regulamentações} emanadas do Conselho Estadual de Educação, deverão promover o credenciamento das \textbf{instituições} de ensino interessadas em \textbf{ofertar} trilhas de aprofundamento.

% matched lemmas: atender, carga, catálogo, certificação, cursar, curso, destinar, diretriz, guia, habilitação, horário, instituição, itinerário, legislação, município, obrigatoriamente, oferta, ofertar, prever, qualificação, regulamentação, técnico, vocação
\end{textsample}
